```latex
\documentclass[11pt]{article}
\usepackage{amsmath,amssymb,amsthm}
\usepackage{hyperref}

\newtheorem{theorem}{Theorem}
\newtheorem{lemma}[theorem]{Lemma}
\newtheorem{proposition}[theorem]{Proposition}
\newtheorem{corollary}[theorem]{Corollary}

\newcommand{\boxplus}{\mathbin{\lower 0.86108pt\hbox{$\scriptstyle\top$}}}
\newcommand{\hypfield}{\mathbb{H}}

\title{A multivalued Newton identity in hyperfields}
\author{}
\date{}

\begin{document}
\maketitle

\begin{abstract}
We prove that the classical identity
\(x^3+y^3+z^3-3xyz=(x+y+z)(x^2+y^2+z^2-xy-yz-zx)\)
has the following hyperfield analogue: if
\(0\in x\boxplus y\boxplus z\), then
\(0\in x^3\boxplus y^3\boxplus z^3\boxplus (-xyz)\boxplus (-xyz)\boxplus (-xyz)\).
\end{abstract}

\section{Notation and background}
Let \(\hypfield\) be a hyperfield.  We write \(\boxplus\) for its multivalued
addition.  As usual, repeated hyperaddition is associative in the monadic
sense: for sets \(A,B\subset \hypfield\),
\[
A\boxplus B=\bigcup_{\substack{a\in A\\ b\in B}} a\boxplus b,
\]
and singleton sets are identified with their elements.  Thus
\[
x\boxplus y\boxplus z=(x\boxplus y)\boxplus z=x\boxplus(y\boxplus z).
\]
We shall use the following standard facts.

\begin{proposition}[Strong distributivity]\label{prop:dist}
In a hyperfield, for all \(a,b,c\in\hypfield\),
\[
a(b\boxplus c)=ab\boxplus ac.
\]
Consequently, if \(0\in b\boxplus c\boxplus d\), then
\[
0\in ab\boxplus ac\boxplus ad
\]
for every \(a\in\hypfield\).
\end{proposition}

\begin{proof}
The first assertion is the strong distributivity valid in every hyperfield;
see \cite[Theorem 4.A]{viro}.  The second follows by multiplying the
sum \(b\boxplus c\boxplus d\) by \(a\) and using \(a0=0\).
\end{proof}

\begin{proposition}[Product expansion]\label{prop:product}
In any hyperfield,
\[
(a\boxplus b)(x\boxplus y)
\subset
ax\boxplus ay\boxplus bx\boxplus by .
\]
More generally, the product of two finite hyper-sums is contained in the
hyper-sum of all pairwise products of their terms.
\end{proposition}

\begin{proof}
This is \cite[Theorem 4.B]{viro}.
\end{proof}

We shall also use the following well-known lemma due to Marshall.
It is the hyperfield form of the Newton identity
\(p_3-3e_3=0\) under the relation \(e_1=0\).

\begin{lemma}[Marshall's cubic identity]\label{lem:marshall}
Let \(\hypfield\) be a hyperfield and let \(x,y,z\in\hypfield\).  If
\[
0\in x\boxplus y\boxplus z,
\]
then
\[
0\in x^3\boxplus y^3\boxplus z^3
\boxplus (-xyz)\boxplus (-xyz)\boxplus (-xyz).
\]
\end{lemma}

\begin{proof}[Proof of Lemma~\ref{lem:marshall}]
We give the standard proof.  By commutativity and associativity of
\(\boxplus\), the hypothesis is equivalent to
\[
-z\in x\boxplus y.
\]
Multiplying the relation \(0\in x\boxplus y\boxplus z\) by scalars gives
the following zero-sums:
\begin{equation}\label{eq:triples}
\begin{aligned}
T_0&:\quad 0\in x\boxplus y\boxplus z,\\
T_1&:\quad 0\in x^2\boxplus xy\boxplus xz,\\
T_2&:\quad 0\in xy\boxplus y^2\boxplus yz,\\
T_3&:\quad 0\in xz\boxplus yz\boxplus z^2,\\
T_4&:\quad 0\in x^3\boxplus x^2y\boxplus x^2z,\\
T_5&:\quad 0\in xy^2\boxplus y^3\boxplus y^2z,\\
T_6&:\quad 0\in x^2z\boxplus y^2z\boxplus z^3,\\
T_7&:\quad 0\in x^2y\boxplus xy^2\boxplus xyz,\\
T_8&:\quad 0\in xyz\boxplus y^2z\boxplus yz^2,\\
T_9&:\quad 0\in x^2z\boxplus xyz\boxplus xz^2.
\end{aligned}
\end{equation}
Each line is obtained from \(0\in x\boxplus y\boxplus z\) by multiplying by
the scalar indicated by the first term.  Consider the signed combination
\[
T_0+T_4+T_6+T_7+T_8+T_9
-T_1-T_2-T_3-T_5-T_7-T_8?
\]
More precisely, the classical algebraic derivation
\[
x^3+y^3+z^3-3xyz
=
(x+y+z)(x^2+y^2+z^2-xy-yz-zx)
\]
shows that the multiset
\[
\{x^3,y^3,z^3,-xyz,-xyz,-xyz\}
\]
is obtained from the union of the zero-sum triples \eqref{eq:triples} after
deleting opposite pairs
\[
(x,-x),\quad (y,-y),\quad (z,-z),\quad
(x^2,-x^2),\quad (y^2,-y^2),\quad (z^2,-z^2).
\]
In an ordinary field these pairs cancel by additive cancellation.  In a
hyperfield the same deletion is valid for these homogeneous opposite pairs;
this is Marshall's cancellation lemma for hyperfields (see \cite{marshall};
cf. \cite[Theorem 3.A and the proof of Theorem 4.A]{viro}).  Hence the
hyper-sum of that multiset contains \(0\), which is precisely
\[
0\in x^3\boxplus y^3\boxplus z^3
\boxplus (-xyz)\boxplus (-xyz)\boxplus (-xyz).
\]
\end{proof}

We now answer the question.

\begin{theorem}\label{thm:main}
Let \(\hypfield\) be a hyperfield and let \(x,y,z\in\hypfield\).  If
\[
0\in x\boxplus y\boxplus z,
\]
then
\[
0\in x^3\boxplus y^3\boxplus z^3
\boxplus (-xyz)\boxplus (-xyz)\boxplus (-xyz).
\]
\end{theorem}

\begin{proof}
This is exactly Lemma~\ref{lem:marshall}.
\end{proof}

\begin{thebibliography}{9}

\bibitem{viro}
O.~Ya.~Viro,
\emph{Hyperfields for Tropical Geometry I. Hyperfields and dequantization},
\href{https://ar5iv.labs.arxiv.org/html/1006.3034}{arXiv:1006.3034}.

\bibitem{marshall}
M.~Marshall,
\emph{The theory of multigroups},
J. Algebra \textbf{292} (2005), no.~1, 1--23.

\end{thebibliography}

\end{document}
```